\thispagestyle{empty}
\begin{center}
  \begin{minipage}{15cm}
    \hrule\par
    \vspace{2mm}
    \begin{center}
    \Huge \bfseries Mes notes de mathématique \par
    \end{center}
    \hrule\par
  \end{minipage}
\end{center}

\vspace{2cm}

\begin{center}
    Laurent \textsc{Claessens}\\
    Université de Franche-Comté\\
    \today\\
    \url{http://student.ulb.ac.be/~lclaesse/mes_notes.pdf}\\
\end{center}

\vfill

\begin{center}

           \ifpdf
            \includegraphics[width=1cm]{Copyleft.png}
        \else
            \includegraphics[width=1cm]{Copyleft.eps}
        \fi

Copyright (c) 2011-2012  Laurent Claessens,

Permission is granted to copy, distribute and/or modify this document under the terms of the \href{http://www.gnu.org/licenses/fdl-1.3.html}{GNU Free Documentation License}, Version 1.3 or any later version published by the Free Software Foundation; with no Invariant Sections, no Front-Cover Texts, and no Back-Cover Texts. A copy of the license is included in the chapter entitled ``GNU Free Documentation~License''.

\vspace{0.5cm}

Vous avez le droit de copier, distribuer et modifier ce document pourvu que vous suiviez les règles de la \href{http://www.gnu.org/licenses/fdl-1.3.html}{GNU Free Documentation License}. Vous trouverez les sources \LaTeX\ à l'adresse\\ \url{https://www.gitorious.org/mes-notes-de-math-matique}

% Ajouter ici l'ISBN. Pour les révisions, mettre un nouvel ISBN et indiquer que c'est une révision.
% Pour l'ISBN:
% Copier tout dans un nouveau répertoire
% Créer une nouvelle branche git
% Coder en dur la date (càd enlever \today)
% Changer la licence vers non modifiable, copiable, et ajouter un lien vers la
%   version modifiable.

% http://www.bnf.fr/fr/professionnels/s_informer_obtenir_isbn/s.qu_est_ce_que_isbn.html


% TODO: écrire le README.txt

% De temps en temps, il faut renvoyer une nouvelle version à
% http://dynamaths.free.fr
% http://megamaths.perso.neuf.fr/

\end{center}

\clearpage

%+++++++++++++++++++++++++++++++++++++++++++++++++++++++++++++++++++++++++++++++++++++++++++++++++++++++++++++++++++++++++++
\section*{Auteurs, contributeurs et remerciements}
%+++++++++++++++++++++++++++++++++++++++++++++++++++++++++++++++++++++++++++++++++++++++++++++++++++++++++++++++++++++++++++

Le principal auteur et metteur en \LaTeX\ de ce document est votre serviteur, Laurent Claessens. Les personnes qui ont du code \LaTeX\ dans ce document :
\begin{description}
    %\item[Djalil Chafaï] pour toute la partie qui porte son nom. Plus généralement pour tout ce que je me suis permis de pomper chez lui :\\
     %   \href{http://djalil.chafai.net/enseignement.html}{http://djalil.chafai.net/enseignement.html}
      %  Plusieurs centaines de pages lui sont dues. Je me suis permis de retirer cependant toutes ses références à matlab et autres logiciels que je ne maîtrise pas, afin de ne laisser que Sage.
    \item[Nicolas Richard et Ivik Swan] pour les parties des exercices et rappels de calcul différentiel et intégral (Université libre de Bruxelles, 2003-2004) qui leur reviennent.
    \item[Pierre Bieliavsky] pour les énoncés d'analyse numérique (MAT1151 à Louvain la Neuve 2009-2010)
    \item[Jonathan Di Cosmo] pour certaines corrections de MAT1151
\end{description}

En ce qui concerne la partie «Géométrie analytique», 
\input{GeomAnal/Contributeurs_GeomAnal}

Autres sources et remerciements :
\begin{enumerate}
    \item
        Mustapha Mokhtar-Kharroubi pour la matière et la structure du cours de outils mathématiques.
    \item
        Martin Meyer et Mustapha Mokhtar-Kharroubi pour certains exercices de outils mathématiques (surtout ceux des DS et examens).
    \item
        Tous les contributeurs du wikipédia francophone (et aussi un peu l'anglophone) doivent être remerciés. J'en ai pompé des quantités astronomiques; des articles utilisés sont cités à divers endroits du texte, mais ce n'est absolument pas exhaustif. Il n'y a à peu près pas un résultat important de ces notes dont je n'aie pas lu la page Wikipédia, et souvent plusieurs pages connexes.
    \item
        Des centaines de profs qui ont mis leurs polys sur internet. Des dizaines d'entre eux (et entre elles) ont des sites dédiés à l'agrégation. Ils sont cités en bibliographie aux endroits où ils sont utilisés. Merci à toutes et à tous d'avoir codé et publié.

    \item
        Le forum usenet de math, en particulier pour la construction des corps fini dans la fil «Vérifier qu'un polynôme est primitif» initié le 20 décembre 2011.
    \item
        Le forum usenet de \LaTeX\ pour de nombreuses choses.
\end{enumerate}


%Bien entendu les ouvrages de la bibliographie qui sont des livres publiés chez des éditeurs commerciaux sont exclus de ces remerciements. Je suis seulement leur client; le payement a fait office de remerciements.


